%% Criado por Daniel Mognato Vaneli - 2019-11-14

\documentclass[12pt,openright,oneside,a4paper,brazil]{abntex2} %twoside imprime uma pagina a mais em branco pro caso de frente e verso
\usepackage[alf]{abntex2cite} % Citações padrão ABNT
\usepackage{lipsum,mathptmx,etoolbox} %semelhante ao Times New Roman
\usepackage[utf8]{inputenc} %acentos
\usepackage{graphicx} %imagens
\usepackage[bottom=2.5cm,top=2.5cm,left=3.0cm,right=3.0cm]{geometry} %geometria da página
\usepackage{indentfirst} %indenta primeiro parágrafo


\usepackage{hyperref} % controla a formação do índice
\titulo{COLOQUE SEU TÍTULO AQUI}
\autor{SEU NOME AQUI}
\local{LOCAL}
\data{ANO}
\orientador{SEU ORIENTADOR}
\instituicao{Universidade Federal Fluminense}
\tipotrabalho{Dissertação (Mestrado)}
\preambulo{PRÉAMBULO AQUI}
\hypersetup{colorlinks,citecolor=black,filecolor=black,linkcolor=black,urlcolor=black}


\makeatletter
\hypersetup{
pdftitle={\@title},
pdfauthor={\@author},
pdfsubject={\imprimirpreambulo}
}
\makeatother

% ----------------------------------------------------------
% Começo do documento
% ----------------------------------------------------------
\begin{document}

% ----------------------------------------------------------
% CAPA
% ----------------------------------------------------------
\begin{figure}[h]
    \flushleft
    \includegraphics[scale=3.0]{imagens/PGMEC2.eps}
\end{figure}
        
    \vspace{0.3cm}
        
    \noindent \large{PROGRAMA FRANCISCO EDUARDO MOURÃO SABOYA DE\\
    PÓS-GRADUAÇÃO EM ENGENHARIA MECÂNICA\\
    ESCOLA DE ENGENHARIA\\
    UNIVERSIDADE FEDERAL FLUMINENSE}
\begin{center}
    \vspace{50pt}
    \Large DISSERTAÇÃO DE MESTRADO\\
    \vspace{100pt}
    \LARGE \imprimirtitulo\\
    \vspace{100pt}
    \Large \imprimirautor\\
    \vspace{\fill}
    \large DEZEMBRO DE \imprimirdata
\end{center}

% ----------------------------------------------------------
% FOLHA DE ROSTO
% ----------------------------------------------------------
\newpage
\begin{center}
    \Large \imprimirautor\\
    \vspace{2cm}
    \LARGE \imprimirtitulo\\
    \vspace{2.5cm}
    \hspace{.45\textwidth}
    \begin{minipage}{.5\textwidth}
        \large{Dissertação de Mestrado apresentada ao Programa Francisco Eduardo Mourão Saboya de Pós-Graduação em Engenharia Mecânica da Universidade Federal Fluminense como parte dos requisitos para a obtenção do título de Mestre em Ciências em Engenharia Mecânica.}
    \end{minipage}%
    \vspace{2.5cm}
    \Large Orientador: \imprimirorientador, D.Sc.\\
    \vspace{\fill}
    \large \imprimirinstituicao\\
    \large \imprimirlocal\\
    \large \imprimirdata
\end{center}

% ----------------------------------------------------------
% FOLHA DE ASSINATURAS
% ----------------------------------------------------------
\newpage
\begin{center}
    \LARGE \imprimirtitulo\\
    \vspace{2.5cm}
    \large Esta dissertação é parte dos pré-requisitos para a obtenção do título de\\
    \vspace{0.2cm}
    \Large \textbf MESTRE EM ENGENHARIA MECÂNICA\\
    \vspace{0.5cm}
    \large Área de concentração: SUA ÁREA DE CONCENTRAÇÃO\\
    \vspace{0.8cm}
    \large Aprovada em sua forma final pela Banca Examinadora formada pelos professores:
    \vspace{2cm}
    \assinatura{\imprimirorientador, D.Sc. \\ Orientador}
    \vspace{1cm}
    \assinatura{Professor \\ Convidado 1}
    \vspace{1cm}
    \assinatura{Professor \\ Convidado 2}
\end{center}

% ----------------------------------------------------------
% DEDICATÓRIA
% ----------------------------------------------------------
\newpage
\begin{dedicatoria}
   \vspace*{\fill}
   %\centering
   \noindent
   \hfill \textit{dedicatória aqui.}
   %\vspace*{\fill}
\end{dedicatoria}

% ----------------------------------------------------------
% AGRADECIMENTOS
% ----------------------------------------------------------
\begin{agradecimentos}

AGRADECIMENTOS AQUI

\end{agradecimentos}

% ----------------------------------------------------------
% EPÍGRAFE
% ----------------------------------------------------------
\begin{epigrafe}
    \vspace*{\fill}
	\begin{flushright}
		\textit{epigrafe aqui\\
		autor da epigrafe}
	\end{flushright}
\end{epigrafe}

% ----------------------------------------------------------
% RESUMO
% ----------------------------------------------------------
\setlength{\absparsep}{18pt} % ajusta o espaçamento dos parágrafos do resumo
\begin{resumo}
 Segundo \cite{ibge1993}, o ...

 \textbf{Palavras-chave}: a, b, c, d.
\end{resumo}

% ----------------------------------------------------------
% LISTA DE FIGURAS
% ----------------------------------------------------------
\pdfbookmark[0]{\listfigurename}{lof}
\listoffigures*
\cleardoublepage

% ----------------------------------------------------------
% LISTA DE TABELAS
% ----------------------------------------------------------
\pdfbookmark[0]{\listtablename}{lot}
\listoftables*
\cleardoublepage

% ----------------------------------------------------------
% LISTA DE SÍMBOLOS
% ----------------------------------------------------------
\begin{simbolos}
\item[$ \Gamma $] Letra grega Gama
\item[$ \Lambda $] Lambda
\item[$ \zeta $] Letra grega minúscula zeta
\item[$ \in $] Pertence
\end{simbolos}

% ----------------------------------------------------------
% SUMÁRIO
% ----------------------------------------------------------
\pdfbookmark[0]{\contentsname}{toc}
\tableofcontents*
\cleardoublepage

% ----------------------------------------------------------
% ELEMENTOS TEXTUAIS
% ----------------------------------------------------------
\textual

% ----------------------------------------------------------
% Introdução (exemplo de capítulo sem numeração, mas presente no Sumário)
% ----------------------------------------------------------
\chapter{Introdução}
\setcounter{page}{1} %começa a contagem das páginas em 1

% ----------------------------------------------------------
\chapter{Motivação}

% ----------------------------------------------------------
% ELEMENTOS PÓS-TEXTUAIS
% ----------------------------------------------------------
\postextual
% ----------------------------------------------------------

% ----------------------------------------------------------
% Referências bibliográficas
% ----------------------------------------------------------
\bibliography{modelo} %chama o arquivo modelo.bib onde estão as referências

\end{document}